\begin{definition}[Read-once branching program]
A $(n,w)$-read-once branching program (ROBP) $B$ is a directed graph on the vertex set $V=\bigcup_{i=0}^n V_i$, where each set $V_i$ contains $w$ nodes. Every edge in this directed graph is labeled either $0$ or $1$. For every $i<n$, and every node $v\in V_i$, there exists exactly two edges starting from $v$, one with label $0$ and the other with label $1$. Every edge starting from a node in $V_i$ connects to a node in $V_{i+1}$.
We say $n$ is the \emph{length} of $B$, $w$ is the \emph{width} of $B$ and $V_i$ is the $i$-th \emph{layer} of $B$.

Moreover, there exists exactly one starting state $s\in V_0$, and exactly one accepting state $t\in V_n$. For every $x=(x_1,\dots,x_n)\in\bits{n}$, we define $B(x)=1$ if starting from $s$ we will reach $t$ following the edges labeled by $x_1,\dots,x_n$. Otherwise we define $B(x)=0$.
\end{definition}
It is well-known the computation of a probabilistic Turing machine that uses space $S$ and tosses $n$ coins,  on a given input $y$, can be carried out by a $(n,2^{O(S)})$-ROBP $B_y$. In particular, if the string $x \in \zo^n$ corresponds to the $n$ coin tosses, then  $B_y(x)$ is the output of the Turing machine. 

A standard derandomization technique is via \emph{pseudorandom generators}. We define this notion for the class of ROBPs.

\begin{definition}[Pseudorandom generator]
A function $G:\bits{s}\to\bits{n}$ is a $(n,w,\eps)$-pseudorandom generator (PRG) if for every $(n,w)$-ROBP $B$,
$$\abs{\ex[x\in\bits{n}]{B(x)}-\ex[r\in\bits{s}]{B(G(r))}}\le\eps.$$
The \emph{seed length} of $G$ is $s$.   $G$ is explicit if $G$ is computable in $O(s)$ space.
\end{definition}
To derandomimze  space-bounded computation  given an explicit $(n,w,\eps)$-PRG, one can enumerate $B(G(r))$ for every $r\in\bits{s}$ with $O(s)$ additional space to compute an $\eps$-approximation of the quantity $\ex[x]{B(x)}$. 

Nisan~\cite{Nis92} constructed a $(n,w,\eps)$-PRG with seed length $O(\log n \cdot \log (nw/\eps))$, which implies $\BPL\subseteq \L^2$.   While there is a lot of progress in constructing PRG with better seed length for restricted family of ROBP (see, e.g., \cite{NZ96,Arm98,BV10,BDVY13,BRRY14,KNP11,De11,Ste12,MRT19} and references therein), Nisan's generator and its variants~\cite{Nis92,INW94,RR99} remain the best-known generators in the general case.
\subsection{Pseudorandom pseudodistribution }
Recently, a beautiful work of Braverman, Cohen and Garg \cite{BCG18} introduced the notion of a \emph{pseudorandom pseudodistribution} (PRPD)  that relaxes the definition of a PRG.
\begin{definition}[Pseudorandom pseudodistribution]
A pair of functions $(G,\rho):\bits{s}\to \bits{n}\times \bbR$ generates a $(n,w,\eps)$-pseudorandom pseudodistribution (PRPD) if for every $(n,w)$-ROBP $B$,
$$\abs{\ex[x\in\bits{n}]{B(x)}-\ex[r\in\bits{s}]{\rho(r)\cdot B(G(r))}}\le\eps.$$
We say $s$ is the \emph{seed length} of $(G,\rho)$. We say $(G,\rho)$ is $k$-bounded if $|\rho(x)|\le k$ for every $x\in\bits{s}$. We say $(G,\rho)$ is explicit if they are computable in space $O(s)$.
\end{definition}
